\section{Interaction Scaling at a Matched Token Budget}
\label{sec:interaction}

\paragraph{The saturation experiment (P1).} On the hard code suite, at a fixed per-task token budget, we compare four strategies: reasoning-only (extended thinking), best-of-$N$ sampling with an \emph{oracle} verifier, a single-agent loop, and the proposer--reviewer harness. The result is \Cref{fig:scaling}, previewed in \Cref{sec:ceiling}: both internal strategies flatten as the budget grows, while \emph{every} strategy that iterates on feedback climbs past them toward a perfect pass rate. The best-of-$N$ comparison is the sharpest, because its verifier is the ground-truth test suite itself, so its ceiling is effectively pass@$N$: the tasks it never solves are those where none of its samples passes, and a perfect selector cannot return a candidate the proposer never writes. The harness solves those same tasks by conditioning the next draw on execution feedback. The gap is information-theoretic, not a tuning artifact.

\paragraph{Code as the clean case study.} Code is the cleanest demonstration, because the instrument is exact and the score is objective. The harness recovers every first-shot failure with no regressions, on both the development suite and a harder deep-spec suite of from-scratch implementations (a JSON parser, a spreadsheet-reference resolver, a minimal edit-script diff, an expression evaluator), lifting both to a perfect $100\%$ pass rate (\Cref{fig:code}). The mechanism is the same across recovered traces: turn~1 produces code that compiles but mishandles an edge case (an off-by-one boundary, semantic-version comparison, CIDR arithmetic); turn~2 receives the failing assertion, the offending input, and the expected-vs-actual values, and rewrites just that block; turn~3 confirms the fix. Reasoning alone cannot supply this: the model cannot know its boundary condition is wrong until a concrete test exercises it.

\begin{figure}[t]
\centering
\begin{tikzpicture}
\begin{axis}[paperbar,width=0.62\linewidth,height=4.8cm,
  ybar,bar width=15pt,ymin=0,ymax=118,
  ylabel={pass rate (\%)},
  symbolic x coords={dev,deepspec},xtick=data,
  xticklabels={development suite (15 tasks),deep-spec suite (11 tasks)},
  x tick label style={font=\footnotesize,color=cMut},
  enlarge x limits=0.45,ytick={0,25,50,75,100},ymajorgrids,
  legend style={at={(0.5,1.03)},anchor=south,legend columns=2,
    /tikz/every odd column/.append style={column sep=2pt}}]
  \addplot[swatch,draw=none,fill=black!22,
    error bars/.cd,y dir=both,y explicit,error bar style={cInk!60,line width=0.6pt}]
    coordinates {(dev,66.7) +- (0,6.7) (deepspec,69.7) +- (0,0)};
    \addlegendentry{single-shot}
  \addplot[swatch,draw=none,fill=cGnd!80]
    coordinates {(dev,100) (deepspec,100)};
    \addlegendentry{reviewed (execution feedback)}
  % selective labels: the reviewed endpoints and the lifts
  \node[font=\scriptsize\bfseries,cInk,anchor=south] at (axis cs:dev,101) {\hspace{15pt}100.0};
  \node[font=\scriptsize\bfseries,cInk,anchor=south] at (axis cs:deepspec,101) {\hspace{15pt}100.0};
  \node[font=\scriptsize,cMut,anchor=south] at (axis cs:dev,76) {\hspace{-15pt}66.7};
  \node[font=\scriptsize,cMut,anchor=south] at (axis cs:deepspec,71) {\hspace{-15pt}69.7};
\end{axis}
\end{tikzpicture}
\caption{\textbf{Execution feedback recovers every first-shot failure on both code suites.} Single-shot vs.\ harness-reviewed pass rate (development suite: $3$-seed mean, whisker $\pm1$\,SD; deep-spec: sign test $p=0.002$). Every failing run is recovered and no passing run regresses; the reviewed bars carry zero seed variance.}
\label{fig:code}
\end{figure}

\paragraph{Architecture buys efficiency and reliability, not ceiling.} Among interaction strategies, the ceiling pass rate ties within seed noise; what separates them is cost and variance (\Cref{fig:efficiency}). The proposer--reviewer harness is the most token-efficient, and the only variant that reaches $100\%$ in \emph{every} seed. Three things explain the efficiency: the reviewer sees only what it needs to locate defects, it returns a structured list rather than raw \texttt{stderr}, and the two roles are specialized.

\begin{figure}[t]
\centering
\begin{tikzpicture}
\begin{axis}[paperbar,width=0.52\linewidth,height=3.6cm,
  xbar,xmin=0,xmax=1600,bar width=11pt,enlarge y limits=0.3,
  xlabel={mean output tokens per task at $B{=}20$K},
  symbolic y coords={loop,iad,harness},ytick={loop,iad,harness},
  yticklabels={single-agent loop,sample-and-select (IAD),proposer--reviewer},
  y tick label style={font=\footnotesize,color=cInk},
  xtick={0,500,1000,1500},xticklabels={0,500,1000,1500},xmajorgrids,clip=false]
  \addplot[draw=none,fill=cGnd!80] coordinates {(1029,harness)};
  \addplot[draw=none,fill=cGnd!45] coordinates {(1416,iad) (1431,loop)};
  % token value at each bar tip; reliability as a right-hand column
  \node[font=\scriptsize\bfseries,cInk,anchor=west,inner sep=2pt] at (axis cs:1029,harness) {1{,}029};
  \node[font=\scriptsize,cMut,anchor=west,inner sep=2pt,fill=white] at (axis cs:1416,iad) {1{,}416};
  \node[font=\scriptsize,cMut,anchor=west,inner sep=2pt,fill=white] at (axis cs:1431,loop) {1{,}431};
  \node[font=\scriptsize\bfseries,cGnd,anchor=west] at (axis cs:1680,harness) {100\% in 3/3 seeds};
  \node[font=\scriptsize,cMut,anchor=west] at (axis cs:1680,iad) {100\% (1 seed)};
  \node[font=\scriptsize,cMut,anchor=west] at (axis cs:1680,loop) {100\% in 2/3 seeds};
\end{axis}
\end{tikzpicture}
\caption{\textbf{Same ceiling, different cost and reliability.} All three interaction strategies use execution feedback and tie on pass rate within seed noise (sign test $p>0.6$); the proposer--reviewer harness converges with ${\sim}28\%$ fewer tokens and is the only variant at $100\%$ across all seeds (\Cref{tab:LvsH}).}
\label{fig:efficiency}
\end{figure}

\paragraph{Built-in early stopping.} The loop does not over-revise: on code, \emph{every} already-passing single-shot task is submitted at iteration~1 with no revision and no wasted tokens. The grounded signal (a failing test) is the trigger; without it the loop is silent. This is why the harness's token overhead over single-shot stays modest despite a generous iteration cap.
